\documentclass[letterpaper,10pt]{article}
\usepackage[utf8]{inputenc}
\usepackage[T1]{fontenc}
\usepackage{lmodern}
\usepackage{microtype}
\usepackage[margin=0.45in]{geometry}
\usepackage{hyperref}
\usepackage{titlesec}
\usepackage{enumitem}
\usepackage{xcolor}

\hypersetup{
    colorlinks=true,
    linkcolor=blue,
    urlcolor=blue,
    pdftitle={Saul Cordoba - CV Fullstack DevOps},
    pdfauthor={Saul Cordoba},
    pdfsubject={Fullstack Engineer \& DevOps | .NET, Python, Angular/React, Docker, CI/CD},
}

\pagestyle{empty}
\urlstyle{same}

\titleformat{\section}{\large\bfseries}{}{0em}{}[\titlerule]
\titlespacing*{\section}{0pt}{5pt}{3pt}

\setlist[itemize]{leftmargin=1.2em, itemsep=1pt, topsep=1.5pt, parsep=0pt, partopsep=0pt}

\begin{document}

%--- HEADER ---
\begin{center}
    {\Huge \bfseries SAÚL CÓRDOBA} \\[3pt]
    {\Large \textbf{Fullstack Engineer \& DevOps | .NET, Python, Angular/React, Docker, CI/CD}} \\[2pt]
    Villavicencio, Colombia (Disponible para Trabajo Remoto / Relocalización) \\
    (+57) 321 961 1182 | \href{mailto:sauleduardo1111@gmail.com}{sauleduardo1111@gmail.com} | \href{https://github.com/Garetetete}{GitHub: Garetetete} | \href{https://www.linkedin.com/in/saul-cordoba-a2997427b}{LinkedIn} | \textbf{Inglés:} B2
\end{center}

%--- PERFIL PROFESIONAL ---
\section{PERFIL PROFESIONAL}
\textbf{Desarrollador Fullstack con perfil DevOps}: construyo el producto y opero su infraestructura. \textbf{Diseñé, desarrollé y llevé a producción desde cero, en solitario y de extremo a extremo, una plataforma SaaS multi-tenant} con backend \textbf{FastAPI/Python} y microservicio \textbf{C\#/.NET}, 120+ endpoints REST/SSE y PostgreSQL con seguridad a nivel de fila — orquestada en \textbf{8 servicios Docker Compose} con Keycloak, nginx, backups automatizados y entornos QA/producción. Experiencia con \textbf{microservicios .NET + Angular} y pipelines \textbf{CI/CD en Azure DevOps y GitHub Actions}. Ingeniero Electrónico con doble especialización (Instrumentación y Control Industrial; Ingeniería de Software).

%--- HABILIDADES TÉCNICAS ---
\section{HABILIDADES TÉCNICAS}
\begin{itemize}
    \item \textbf{Lenguajes:} C\# (.NET), Python, TypeScript/JavaScript, Java, C, C++, SQL.
    \item \textbf{Backend \& Frontend:} .NET Core (Clean Architecture), FastAPI, Node/Express, microservicios, API gateways, SSE, Angular, React, TailwindCSS.
    \item \textbf{Contenedores y despliegue:} Docker y Docker Compose multi-entorno (dev/QA/test), healthchecks, autoheal, límites de recursos, nginx (reverse proxy), Linux, SSH/despliegue remoto en servidores, Railway.
    \item \textbf{CI/CD:} Azure DevOps Pipelines, GitHub Actions, despliegue continuo, quality gates (ruff, pytest), scripting PowerShell/Bash.
    \item \textbf{Bases de datos y seguridad:} PostgreSQL (RLS, JSONB, migraciones, backups pg\_dump), Oracle, SQL Server, MongoDB, Redis, Keycloak (OAuth2/OIDC), JWT, cifrado de secretos (Fernet).
    \item \textbf{Metodologías:} SDD (Spec-Driven Development), TDD, Clean Architecture, OpenSpec.
\end{itemize}

%--- EXPERIENCIA ---
\section{EXPERIENCIA PROFESIONAL}

\textbf{Desarrollador de Software} \hfill \textbf{2025 -- Presente} \\
\textit{SASOFTCO} \hfill \textit{Villavicencio, Colombia}
\begin{itemize}
    \item \textbf{Creador de una plataforma SaaS de análisis documental con IA} (\textbf{en producción}): la construí completa desde cero — backend FastAPI (Python 3.11) por capas, \textbf{120+ endpoints REST/SSE}, frontend server-rendered (Jinja2 + JS modular) y microservicio \textbf{C\#/.NET (Clean Architecture)} para conversión de documentos.
    \item \textbf{PostgreSQL 15 multi-tenant con seguridad a nivel de fila (RLS)}: middleware que resuelve el cliente por request, migraciones versionadas y cifrado de secretos en reposo; \textbf{Keycloak (OAuth2/OIDC)} con validación JWT y permisos por vista.
    \item \textbf{Operación de la plataforma}: orquestación de \textbf{8 servicios con Docker Compose} (API, .NET, LibreOffice headless, nginx, Keycloak, PostgreSQL, backups, autoheal), healthchecks dependientes, entornos separados dev/QA/test, backups automáticos cada 24h con retención de 14 días y scripts de operación en PowerShell/Bash.
    \item \textbf{Integraciones}: consumo de APIs de datos abiertos gubernamentales (Socrata/SoQL) con paginación y control de errores, APIs comerciales, almacenamiento configurable en la nube (OneDrive/SharePoint, Google Drive) y notificaciones por Gmail OAuth.
    \item Desarrollo backend y frontend en \textbf{microservicios .NET con Angular y React} para plataformas de trámites ciudadanos y sector asegurador; \textbf{pipelines CI/CD en Azure DevOps y GitHub Actions}; despliegues remotos por SSH en servidores Linux.
\end{itemize}

\textbf{Desarrollador / Analista de Datos} \hfill \textbf{2024 -- 2025} \\
\textit{Btilink SAS} \hfill \textit{Villavicencio, Colombia}
\begin{itemize}
    \item \textbf{Analista de datos para Aera Energy LLC} (compañía petrolera de EE.UU.): consultas \textbf{SQL sobre Oracle} y dashboards en \textbf{Tableau/Power BI} sobre volúmenes de producción, ingresos y salidas, usados por varios equipos para el monitoreo continuo de la operación.
\end{itemize}

%--- PROYECTOS ---
\section{PROYECTOS DESTACADOS}
\begin{itemize}
    \item \textbf{Sistema de ticketing QR} (Laravel 10 + PostgreSQL, dockerizado): tickets firmados, pagos con webhooks, control de aforo y validación anti doble ingreso.
\end{itemize}

%--- EDUCACIÓN ---
\section{EDUCACIÓN}
\textbf{Especialización en Ingeniería de Software} \hfill \textbf{2025 -- 2026} \\
Universidad de los Llanos \\
\textbf{Especialización en Instrumentación y Control Industrial} \hfill \textbf{2024 -- 2025} \\
Universidad de los Llanos --- \textit{Modalidad de grado: Diplomado en Sistemas IoT y PLC} \\
\textbf{Ingeniería Electrónica} \hfill \textbf{2018 -- 2024} \\
Universidad de los Llanos

%--- CERTIFICACIONES (descomentar y completar cuando tengas los datos) ---
% \section{CERTIFICACIONES Y DISTINCIONES}
% \begin{itemize}
%     \item \textbf{Nombre de la certificación} -- Entidad emisora (Año).
%     \item \textbf{Distinción académica:} promedio destacado, beca o reconocimiento.
% \end{itemize}

\end{document}
